\scanpage{296}
- \mspage{2} -

\uline{il existe une suite spectrale de $\Lambda$-modules aboutissant à}
$\operatorname{Ext}_{\Lambda}^{*}(\varinjlim A_{\alpha},B)$, \uline{et dont le premier terme est donné par}
\[
E_2^{pq}=\varprojlim_{\alpha}^{(p)}\operatorname{Ext}_{\Lambda}^{q}(A_{\alpha},B).
\]

Utilisons maintenant cette proposition avec $B=\Lambda$ (régulier) et $A_{\alpha}=\operatorname{Ext}_{\Lambda}^{k}(L_k(M_{\alpha}),\Lambda)$. Nous obtenons $E_2^{pq}=0$, $q<k$;
$E_2^{pk}=\varprojlim^{(p)}L_k(L_k(M_{\alpha}))=\varprojlim^{(p)}L_k(M_{\alpha})$ et $E_2^{pq}=0$, $p>\dim\Lambda-q-2$, $q>k$, en vertu du théorème 1, resp. $\alpha)$, $\eta)$, $\vartheta)$ de $(^{1})$ et à cause de l'hypothèse de récurrence. Donc la suite spectrale dégénère en l'isomorphisme
\[
\varprojlim^{(p)}L_k(M_{\alpha})=\operatorname{Ext}_{\Lambda}^{p+k}(\varinjlim\operatorname{Ext}_{\Lambda}^{k}(L_k(M_{\alpha}),\Lambda),\Lambda)\quad\text{pour }p\geq\dim\Lambda-k;
\]
et le \deleted{\unsure{$2^e$}} deuxième membre est bien zéro pour $p>\dim\Lambda-k$, ce qui achève l'induction et la démonstration du théorème 2.

REMARQUE 1. - Le théorème 2, restreint au cas $\Lambda=\mathbb Z$, généralise un résultat de Milnor ($(^{2})$, Appendix).

REMARQUE 2. - Le théorème 2 n'est pas valable en général si I n'est pas filtrant, ou si les $M_{\alpha}$ ne sont pas de type fini.

En ce qui concerne les $\varprojlim^{(i)}M_{\alpha}$ avec $i\leq\dim\Lambda$, nous n'avons que des résultats fragmentaires. A l'aide des méthodes de $(^{2})$, on démontre facilement

\margin{\uline{\unsure{Non filtré ?}}}
PROPOSITION 2. - \uline{Soit} A \uline{un $\Lambda$-module} ($\Lambda$ commutatif) \uline{et} $(M_{\alpha})$ \uline{un système projectif de $\Lambda$-modules. Alors il existe deux suites spectrales} de $\Lambda$-modules \uline{aboutissant à la même limite} $H^{*}$, \uline{avec comme premiers termes}
\[
E_2^{pq}=\varprojlim^{(p)}\operatorname{Ext}_{\Lambda}^{q}(A,M_{\alpha})\quad\uline{\text{et}}\quad {}'E_2^{pq}=\operatorname{Ext}_{\Lambda}^{p}(A,\varprojlim^{(q)}M_{\alpha}).
\]

Soit maintenant $\Lambda$ régulier, I filtrant et les $M_{\alpha}$ et A de type fini. Alors $E_2^{pq}=0$ pour $p>\dim\Lambda-q$ [théorème 1, $\delta)$ de $(^{1})$ et théorème 2], d'où $H^s=0$ pour $s>\dim\Lambda$. La deuxième suite spectrale donne $\operatorname{Ext}_{\Lambda}^{1}(A,\varprojlim^{(1)}M_{\alpha})=0$, si $\dim\Lambda=1$, et $\operatorname{Ext}_{\Lambda}^{\nu-1}(A,\varprojlim^{(\nu)}M_{\alpha})=0$, si $\nu=\dim\Lambda>1$, pour tout $\Lambda$-module A de \uline{type fini}. Mais ceci
\margin{\figure{\begin{tikzpicture}\draw (0,0.12)--(0.04,0)--(0.14,0.24);\end{tikzpicture}}}
\mspage{295}
